\chapter[Evaluación detallada]{Evaluación detallada de transporte y formato}
\label{anx:evaluacion}
\small

\section{Pros y contras de cada mecanismo de transporte}

\textbf{MQTT (4,85).} A favor: publicar/suscribir desacopla emisores de consumidores; QoS en tres
niveles sin lógica propia; el testamento (\emph{Last Will}) hace inmediata la detección de caída;
cabecera de 2 bytes; conexión saliente persistente que atraviesa NAT; nació para IoT en redes
inestables. En contra: exige operar un broker (Mosquitto: un contenedor, consumo mínimo); el
navegador no lo habla nativo (irrelevante: el panel consulta la API).

\textbf{HTTP/REST (3,45).} A favor: universal; sin broker; trivial de probar. En contra: cientos de
bytes de cabeceras por decenas de bytes de dato; sin noción de «cliente conectado» (la caída solo se
infiere por silencio); una conexión por envío. Rol: vía secundaria de integración.

\textbf{CoAP (3,05).} El «HTTP compacto» sobre UDP para nodos mínimos. Digno técnicamente;
ecosistema pequeño, UDP a veces filtrado, sin ventaja sobre MQTT en este caso.

\textbf{WebSocket (3,05).} Canal persistente pensado para navegadores; para sensores obliga a
construir encima temas, confirmaciones y testamento: lo que MQTT trae de fábrica. Útil como túnel de
MQTT para paneles en vivo.

\textbf{Socket TCP propio (2,60).} Máxima ligereza y control; a cambio, reinventar sesiones,
reconexión, confirmaciones y formato. Se descarta con argumento: meses de trabajo para un resultado
inferior a MQTT.

\section{Identificación de dispositivos}

El registro (aprovisionamiento) declara identidad, tipo, ubicación y credencial; cada trama se
autoidentifica con el campo \texttt{bn}. La validación cruza ambos: dispositivo desconocido se
rechaza y se registra el intento; trama malformada de dispositivo conocido va a cuarentena; unidad
que no corresponde al tipo declarado va a cuarentena. Un sensor nuevo no exige tocar código: se
registra y sus tramas ya se entienden.

\section{Arquitecturas de nodo consideradas (lado del compañero)}

\begin{longtable}{@{}L{3.4cm}L{3.2cm}L{3.2cm}L{3.2cm}@{}}
\filaencab \encab{} & \encab{Puente ESP32} & \encab{Mini-PC de campo} & \encab{Sensor con red propia} \\ \midrule \endhead
Costo por punto & $\sim$5--10 USD & $\sim$80--200 USD & según fabricante \\
Consumo & bajo, admite batería & alto, enchufado & variable \\
Flexibilidad & cualquier sensor serial & máxima (Python, SSH) & la que dio el fabricante \\
Rol propuesto & nodo estándar & banco de pruebas & caso de integración HTTP \\
\bottomrule
\caption{Comparación de arquitecturas de nodo. La decisión es del trabajo coordinado del compañero;
para esta plataforma todas convergen en el tópico MQTT.}
\end{longtable}

\section{Evaluación de la tecnología de backend}

Cinco alternativas sobre siete criterios. El peso mayor es la experiencia previa del desarrollador
(25\,\%), porque en un proyecto individual con plazos cerrados el tiempo de aprendizaje se descuenta
del de construcción; le siguen la correspondencia entre la estructura del código y la arquitectura
documentable (15\,\%, la Etapa 2 asigna 55\,\% al diseño), la generación automática de la
especificación de la API (15\,\%), la unificación de lenguaje con la interfaz (15\,\%), la madurez de
bibliotecas MQTT/PostgreSQL (15\,\%), el tipado estático (10\,\%) y el rendimiento (5\,\%: todas las
alternativas superan con holgura los 50 mensajes por segundo requeridos).

\begin{table}[H]
\centering\footnotesize
\begin{tablaseg}
\begin{tabular}{@{}L{3.6cm}C{1.3cm}C{1.4cm}C{1.4cm}C{1.4cm}C{1.1cm}C{1.1cm}C{1.1cm}C{1.1cm}@{}}
\toprule
\textbf{Alternativa} & \textbf{Exper.} & \textbf{Arquit.} & \textbf{OpenAPI} & \textbf{1 leng.} & \textbf{Libs} & \textbf{Tipos} & \textbf{Rend.} & \textbf{Total} \\
 & 25\,\% & 15\,\% & 15\,\% & 15\,\% & 15\,\% & 10\,\% & 5\,\% & \\
\midrule
\textbf{TypeScript + NestJS} & 4 & 5 & 5 & 5 & 5 & 5 & 4 & \textbf{4,70} \\
TypeScript + Fastify & 5 & 3 & 4 & 5 & 5 & 5 & 5 & 4,55 \\
Python + FastAPI & 3 & 3 & 5 & 2 & 5 & 4 & 3 & 3,55 \\
Java + Spring Boot & 2 & 5 & 4 & 1 & 5 & 5 & 4 & 3,45 \\
Go + Echo & 1 & 3 & 3 & 1 & 4 & 5 & 5 & 2,65 \\
\bottomrule
\end{tabular}
\end{tablaseg}
\caption{Matriz de tecnología de backend (escala 1--5).}
\end{table}

Se selecciona \textbf{NestJS}: sus módulos corresponden uno a uno con los módulos M1--M8 del
análisis (los diagramas de la Etapa 2 describirán el código real), y con TypeScript en backend e
interfaz los tipos de la trama se definen una sola vez para todo el sistema. La diferencia con
Fastify (0,15) no es decisiva y migrar sería una desviación aceptable si la curva de NestJS
resultara costosa. \textbf{La interfaz (frontend)} se decide en la Etapa 2, con el diseño de vistas
en mano y tres restricciones ya fijadas: TypeScript, componentes de código abierto y uso en
navegador móvil; se evaluarán React, Vue y Svelte con una matriz análoga.

\section{Evaluación de la tecnología de interfaz (frontend)}

El panel es el único módulo que el usuario final ve, pero \textbf{la restricción dominante de este
trabajo es la ligereza}: la plataforma debe operar en hardware modesto, por debajo de las
exigencias del referente (ThingsBoard recomienda 4\,GB de memoria solo para sí mismo). Por ello el
criterio de mayor peso es el tamaño del artefacto entregado al navegador y el consumo en el
servidor que lo sirve, seguido de la coherencia con el backend (TypeScript), la disponibilidad de
componentes de gráficos de series de tiempo, y la experiencia previa.

\begin{table}[H]
\centering\footnotesize
\begin{tablaseg}
\begin{tabular}{@{}L{3.4cm}C{1.6cm}C{1.6cm}C{1.7cm}C{1.5cm}C{1.5cm}C{1.1cm}@{}}
\toprule
\textbf{Alternativa} & \textbf{Ligereza} & \textbf{TypeScript} & \textbf{Gráficos series} & \textbf{Experien.} & \textbf{Ecosist.} & \textbf{Total} \\
 & 35\,\% & 20\,\% & 20\,\% & 15\,\% & 10\,\% & \\
\midrule
\textbf{Svelte + uPlot} & 5 & 5 & 5 & 3 & 3 & \textbf{4,45} \\
Vue + ECharts & 3 & 5 & 4 & 4 & 4 & 3,75 \\
React + Recharts & 2 & 5 & 5 & 4 & 5 & 3,70 \\
Angular (como ThingsBoard) & 1 & 5 & 4 & 2 & 4 & 2,85 \\
\bottomrule
\end{tabular}
\end{tablaseg}
\caption{Matriz de tecnología de interfaz (escala 1--5).}
\end{table}

Se recomienda \textbf{Svelte con uPlot}. Svelte compila a JavaScript mínimo sin \emph{virtual DOM},
lo que produce artefactos varias veces más livianos que React o Angular —coherente con la ligereza
que diferencia a la plataforma de ThingsBoard, construido en Angular—; uPlot es una biblioteca de
gráficos especializada en series de tiempo de unos pocos kilobytes, apropiada para paneles de
telemetría en navegador móvil. Se mantiene TypeScript, con lo que los tipos del modelo interno se
comparten entre backend e interfaz. La decisión se confirma en la Etapa 2 con el diseño de las
vistas en mano; React queda como alternativa si se prioriza el ecosistema de componentes sobre la
ligereza.

\section{Presupuesto de recursos frente al referente}

La preocupación por el consumo de hardware es legítima y se aborda con una estimación explícita del
uso de memoria en reposo de cada componente:

\begin{table}[H]
\centering\footnotesize
\begin{tabular}{@{}L{5.0cm}L{3.2cm}L{4.4cm}@{}}
\toprule
\textbf{Componente} & \textbf{Memoria en reposo} & \textbf{Nota} \\
\midrule
Mosquitto (broker) & $\sim$10\,MB & escrito en C \\
TimescaleDB (PostgreSQL) & $\sim$150--300\,MB & crece con la carga; configurable \\
Backend NestJS (Node) & $\sim$80--150\,MB & \\
ChirpStack (si hay rama LoRa) & $\sim$50\,MB & escrito en Go \\
Interfaz servida (estática) & $\sim$10\,MB & Svelte compilado \\
\midrule
\textbf{Total estimado en reposo} & $\sim$\textbf{0,3--0,5\,GB} & frente a $\geq$4\,GB de ThingsBoard \\
\bottomrule
\end{tabular}
\caption{Presupuesto de memoria estimado. Da origen al requerimiento RNF-17: la plataforma completa
debe operar en un servidor de 2\,GB de memoria, lo que se verifica empíricamente en la Etapa 4.}
\end{table}

Esta estimación es preliminar y se somete a medición real en el banco de comparación de la Etapa 4:
es una hipótesis de diseño, no una afirmación cerrada. El riesgo de que TimescaleDB crezca bajo
carga se controla con políticas de retención y compresión.

\section{Costos de alojamiento (preliminar; el detalle es OE3)}

Servidor virtual arrendado: 6--20 USD/mes (2 vCPU, 4 GB). Mini-PC propio en la organización:
150--250 USD por única vez. Servidor institucional: sin costo si se dispone. Todos los componentes
son de código abierto sin licencias.
